\section{What the efficacy evidence shows}\label{sec:evidence}
\subsection{The case for CBT}
Against placebo, CBT outperforms on anxiety target symptoms ($g=0.56$ across 41 trials; \citealt{carpenter2018cbt}). Against inactive controls, effects are larger: a meta-analysis of 144 trials found $g=0.75$ (depression) to $0.88$ (social anxiety disorder), with $0.80$ for GAD and $0.81$ for panic disorder \citep{cuijpers2016update}. A 2025 network of 375 trials gave anxiety-spectrum effects between $0.5$ and $1.0$, above $1.0$ for PTSD and specific phobia \citep{cuijpers2025unified}. In social anxiety, individual CBT ranked best of the therapies compared in a network meta-analysis of 101 trials \citep{mayowilson2014sad}, and in GAD only CBT remained significantly better than treatment as usual at 3--12 months in a network of 65 trials \citep{papola2024gad}. Any critique has to explain why CBT works at all.

\subsection{The headline shrinks with better controls and bias adjustment}
\begin{itemize}
\item \textbf{Control type.} More than 80\% of anxiety trials used waitlist controls, and the few that did not showed much smaller effects \citep{cuijpers2016update}; only 17.4\% of trials in that analysis were high quality. In the 375-trial analysis, waitlist-controlled effects exceeded $g=0.94$, while care-as-usual comparisons gave $0.22$--$1.13$ \citep{cuijpers2025unified}.
\item \textbf{Bias.} In depression, psychotherapy effects fell from $g=0.70$ to $g=0.31$ after adjusting for several sources of bias, with only 23\% of studies at low risk of bias \citep{cuijpers2019eysenck}. Of NIH-funded psychotherapy trials, 23.6\% never reached publication; including them reduced the pooled effect by about 25\% \citep{driessen2015publicationbias}. An umbrella review of 102 recent meta-analyses (3{,}782 RCTs) found most target-symptom effects small and pooled psychotherapy effects of $0.34$ against placebo or treatment as usual \citep{leichsenring2022umbrella}.
\item \textbf{Response and remission.} About half of anxiety patients respond or remit \citep{loerinc2015response,springer2018remission}.
\end{itemize}

\subsection{Does the effect shrink over time?}
This is contested, and the evidence is mixed. In depression, one meta-analysis reported a steady decline in CBT effects since its introduction \citep{johnsen2015decline}, but a reanalysis restricted to randomised comparisons judged the decline ``most likely a spurious finding'' \citep{cristea2017notfalling}, and a 409-trial meta-analysis found no reduction by publication year \citep{cuijpers2023cbtdepression}. In anxiety, a meta-regression of 49 controlled studies over 30 years found a flat trend (publication year $B=-0.008$, $p=.24$): effects ``did not increase over the last 3 decades'' \citep{hofmann2025effectsizes}. Placebo-controlled anxiety estimates have been reported at lower values in successive reviews (0.73, 0.56, 0.24), but those come from separate reviews with different case mixes, not from one meta-regression. We therefore do \emph{not} claim a decline over time; the defensible statement is that effects have not grown and that they are sensitive to control condition and bias.
